% -----------------------------------------------------------------------------
\subsection{Problema 4.7. Electrete cilíndrico coaxial}\label{sec:p47}

\begin{problema}[Electrete coaxial]\label{prob:electrete-coaxial}
Sea un electrete limitado por dos superficies cilíndricas coaxiales infinitamente largas de radios $a$ y $b$ ($a<b$). El electrete está
polarizado con un momento dipolar por unidad de volumen $\vP=\dfrac{kr}{2}\,\uv r$, siendo $k$ una constante, $r$ la distancia al eje de
simetría ($a<r<b$) y $\uv r$ el vector unitario en la dirección radial. Calcula el campo electrostático $\vE$ en todo el espacio y el vector
desplazamiento $\vD$ en el material.
\end{problema}

\begin{resolucion}
\idea Mismo esquema que el \cref{prob:electrete-esferico}, con simetría cilíndrica (método M2). Ahora la polarización no es solenoidal y sí hay
carga ligada de volumen. No hay carga libre y $\vP$ es irrotacional (no varía en la dirección perpendicular a $\uv r$), luego el problema es
trivial: $\vD=\vcero$ para todo $\vr$.

\paso{Paso 1: campo} $\vE=-\vP/\eo$ en el material y nulo fuera:
\begin{equation}\label{eq:p47}
  \recuadro{\vE=\vcero\ (r<a),\qquad \vE=-\frac{kr}{2\eo}\,\uv r\ \ (a<r<b),\qquad \vE=\vcero\ (r>b);\qquad \vD=\vcero .}
\end{equation}
\paso{Paso 2: cargas ligadas} En coordenadas cilíndricas, $\rhoP=-\dfrac1r\dfrac{\partial}{\partial r}\Big(r\,\dfrac{kr}{2}\Big)=-k$. En las superficies, con la
normal exterior al electrete,
\[
  \sigmaP(a)=\vP\cdot(-\uv r)=-\frac{ka}{2}\quad(\text{apunta hacia el eje}),\qquad \sigmaP(b)=\vP\cdot\uv r=\frac{kb}{2}.
\]
\paso{Paso 3: comprobación con Gauss para $\vE$} Para un cilindro gaussiano de radio $a<r<b$ y longitud $\ell$,
\[
  \qPenc=-\frac{ka}{2}\,2\pi a\ell-k\,\pi(r^2-a^2)\,\ell=-k\pi r^2\ell
  \quad\Longrightarrow\quad
  E\,2\pi r\ell=-\frac{k\pi r^2\ell}{\eo},\quad \vE=-\frac{kr}{2\eo}\,\uv r .
\]
Para $r>b$, $\qPenc=-k\pi a^2\ell+kb\,\pi b\ell-k\pi(b^2-a^2)\ell=0$: la carga de polarización neta es siempre cero.
\end{resolucion}

\begin{observacion}[Diferencia de potencial]\label{obs:p47}
Aunque fuera del electrete no haya campo, entre sus caras sí hay diferencia de potencial:
$\phi(a)-\phi(b)=\int_a^bE_r\,\dd r=-\dfrac{k}{4\eo}(b^2-a^2)$. La superficie interior está a menor potencial, en acuerdo con la carga ligada
negativa que la recubre (\cref{fig:p47}).
\end{observacion}

% Figura: problema 4.7. Electrete cilíndrico coaxial con P = (k r/2) u_r.
\begin{figure}[t]
\centering
\begin{subfigure}[b]{0.27\textwidth}\centering
\begin{tikzpicture}[scale=0.8]
  % cilindro hueco en perspectiva
  \fill[dielec] (-1.3,0) -- (-1.3,-3.0) arc (180:360:1.3 and 0.38) -- (1.3,0) arc (0:180:1.3 and 0.38);
  \fill[dielec] (-1.3,0) arc (180:360:1.3 and 0.38) -- (1.3,0) arc (0:180:1.3 and 0.38);
  \fill[dielecB] (0,0) ellipse (1.3 and 0.38);
  \fill[fondo] (0,0) ellipse (0.55 and 0.16);
  \draw[borde] (0,0) ellipse (1.3 and 0.38) (0,0) ellipse (0.55 and 0.16);
  \draw[borde] (-1.3,0) -- (-1.3,-3.0) arc (180:360:1.3 and 0.38) -- (1.3,0);
  \draw[texto!50!fondo, dashed, line width=0.4pt] (-0.55,0) -- (-0.55,-3.0) (0.55,0) -- (0.55,-3.0);
  \draw[texto!50!fondo, dashed, line width=0.4pt] (-0.55,-3.0) arc (180:360:0.55 and 0.16);
  \foreach \t in {20,70,...,340}
    \draw[campoP, line width=0.7pt, -{Stealth[length=1.3mm]}] ({0.6*cos(\t)},{0.175*sin(\t)}) -- ({1.22*cos(\t)},{0.355*sin(\t)});
  \draw[cota] (1.65,0) -- (1.65,-3.0) node[midway, right, etiq] {$\ell$};
  \node[etiq] at (-0.25,0.52) {$a$};
  \node[etiq] at (1.0,0.62) {$b$};
\end{tikzpicture}
\caption{Electrete coaxial}
\end{subfigure}\hfill
\begin{subfigure}[b]{0.27\textwidth}\centering
\begin{tikzpicture}[scale=0.8]
  \fill[dielec] (0,0) circle (1.7);
  \fill[fondo] (0,0) circle (0.8);
  \draw[borde] (0,0) circle (1.7) (0,0) circle (0.8);
  \foreach \t in {0,45,...,315}
    \draw[campoP, line width=0.8pt, -{Stealth[length=1.5mm]}] (\t:0.9) -- (\t:{0.9+0.62*1.6/1.7});
  \draw[vecn] (135:0.8) -- (135:0.3) node[right=-1pt, etiq] {$\nn$};
  \draw[vecn] (1.7,0) -- (2.3,0) node[above, etiq] {$\nn$};
  \node[etiq, text=campoP] at (62:2.0) {$\vP$};
  \node[etiq, font=\footnotesize] at (0,-0.38) {$\vE=\vcero$};
\end{tikzpicture}
\caption{Corte transversal}
\end{subfigure}\hfill
\begin{subfigure}[b]{0.42\textwidth}\centering
\begin{tikzpicture}
\begin{axis}[width=5.8cm, height=4.3cm, axis lines=middle, xmin=0, xmax=3.1, ymin=-1.15, ymax=1.15,
  xtick={1,2}, xticklabels={$a$,$b$}, ytick=\empty, xlabel={$r$},
  xlabel style={at={(ticklabel* cs:1)}, anchor=north east}, clip=false,
  legend style={at={(0.5,-0.08)}, anchor=north, legend columns=2, font=\tiny, /tikz/every even column/.append style={column sep=4pt}}]
  \addplot[curvaP, domain=1:2, samples=2] {x/2}; \addlegendentry{$P=kr/2$}
  \addplot[curvaE, domain=1:2, samples=2] {-x/2}; \addlegendentry{$\eo E_r=-kr/2$}
  \addplot[curvaD] coordinates {(0,0.02) (3.1,0.02)}; \addlegendentry{$D_r=0$}
  \addplot[curvaP, forget plot] coordinates {(1,0) (1,0.5)};
  \addplot[curvaP, forget plot] coordinates {(2,0) (2,1)};
  \draw[equipot] (axis cs:1,-1.15) -- (axis cs:1,1.15) (axis cs:2,-1.15) -- (axis cs:2,1.15);
  \node[etiqt, anchor=west] at (axis cs:2.05,1.0) {$\sigmaP(b)=\tfrac{kb}{2}$};
  \node[etiqt, anchor=west] at (axis cs:1.03,-0.22) {$\sigmaP(a)=-\tfrac{ka}{2}$};
  \node[etiqt] at (axis cs:1.32,-1.05) {$\rhoP=-k$};
\end{axis}
\end{tikzpicture}
\caption{Perfiles radiales}
\end{subfigure}
\caption{\Cref{prob:electrete-coaxial}. (a), (b) Electrete limitado por dos cilindros coaxiales de radios $a<b$, con polarización
radial creciente. (c) Hay carga ligada de volumen uniforme, $\rhoP=-k$, y de superficie en ambas caras;
su total por unidad de longitud es nulo. Como no hay carga libre y $\vP$ es irrotacional, $\vD=\vcero$ y
$\vE=-\vP/\eo$ dentro del material; fuera de él, el campo es nulo.}
\label{fig:p47}
\end{figure}


% -----------------------------------------------------------------------------
\subsection{Problema 4.8. Condensador esférico con dos dieléctricos}\label{sec:p48}

\begin{problema}[Dos dieléctricos hemisféricos]\label{prob:dos-hemisferios}
Sean dos esferas conductoras concéntricas de radios $a$ y $b$ separadas por dos dieléctricos de permitividades $\varepsilon_1$ y $\varepsilon_2$ cuya superficie
de contacto es un plano ecuatorial (\cref{fig:p48}). Se le proporciona una carga $Q$ a la esfera interior y se conecta la esfera exterior a potencial
nulo. Demuestra que la expresión para el potencial es la misma en las dos zonas del recinto entre las dos esferas. Calcula dicha solución y la
capacidad del sistema. ¿Cómo se distribuye la carga $Q$ en la esfera interior?
\end{problema}

\begin{resolucion}
\idea Probamos un potencial que solo dependa de $r$ en ambas zonas y comprobamos que satisface todas las condiciones; por unicidad, es la
solución. Después, Gauss para $\vD$ (métodos M3 y M4).

\paso{Paso 1: el mismo potencial en ambas zonas} Busquemos $\phi=\phi(r)$ en $a<r<b$. Entonces $\vE=E(r)\,\uv r$ es radial y, por tanto, \emph{tangente} al
plano ecuatorial que separa los medios. Las condiciones de frontera en ese plano son:
\begin{itemize}
  \item componente tangencial de $\vE$ continua: $E_1(r)=E_2(r)\equiv E(r)$ en cada punto del plano, es decir, la misma función en ambos medios;
  \item componente normal de $\vD$ continua: $\vD_i=\varepsilon_iE\,\uv r$ no tiene componente normal al plano, y $0=0$ se cumple trivialmente.
\end{itemize}
Además, en cada medio $\phi$ debe ser armónico, y las soluciones radiales de $\lap\phi=0$ son $\phi=A/r+B$; las dos esferas son equipotenciales y
$\phi(b)=0$. Una única función $\phi(r)=A/r+B$ satisface todo: por el teorema de unicidad~\cite{T3}, la expresión del potencial es la misma en las
dos zonas.

\paso{Paso 2: $\vD$ y $\vE$ por Gauss} Con una esfera de radio $a<r<b$, la mitad en cada medio,
\[
  \oiint\vD\cdot\dS=\varepsilon_1E\,2\pi r^2+\varepsilon_2E\,2\pi r^2=Q
  \quad\Longrightarrow\quad
  \vE=\frac{Q}{2\pi(\varepsilon_1+\varepsilon_2)\,r^2}\,\uv r .
\]
\paso{Paso 3: potencial y capacidad} Con $\phi(b)=0$,
\[
  \recuadro{\phi(r)=-\int_b^r E\,\dd r'=\frac{Q}{2\pi(\varepsilon_1+\varepsilon_2)}\Big(\frac1r-\frac1b\Big),\qquad
  C=\frac{Q}{\phi(a)}=2\pi(\varepsilon_1+\varepsilon_2)\,\frac{ab}{b-a}.}
\]
\paso{Paso 4: reparto de la carga} En la superficie del conductor interior, $\sigmaf=D_n=\varepsilon_iE(a)$ en cada hemisferio:
\[
  \recuadro{\sigma_1=\frac{\varepsilon_1\,Q}{2\pi a^2(\varepsilon_1+\varepsilon_2)},\qquad \sigma_2=\frac{\varepsilon_2\,Q}{2\pi a^2(\varepsilon_1+\varepsilon_2)} .}
\]
Comprobación: $2\pi a^2(\sigma_1+\sigma_2)=Q$.
\end{resolucion}

\begin{observacion}[Dos condensadores en paralelo y la carga total]\label{obs:p48}
La capacidad es la suma $C=C_1+C_2$, con $C_i=2\pi\varepsilon_i\,ab/(b-a)$ la de medio condensador esférico lleno del dieléctrico $i$: los dos hemisferios
están al mismo potencial y se reparten la carga, como dos condensadores en paralelo. La carga libre se concentra en el hemisferio de mayor
permitividad, pero la carga ligada adyacente, $\sigma_{\mathrm{P},i}=-(\varepsilon_i-\eo)E(a)$, compensa justo esa diferencia: la carga total,
$\sigma_i+\sigma_{\mathrm{P},i}=\eo E(a)=Q\eo/[2\pi a^2(\varepsilon_1+\varepsilon_2)]$, es uniforme. Por eso el campo es radial y con simetría esférica, aunque el sistema no la tenga
(\cref{fig:p48}\,(b)). El problema se resolvió en los seminarios con exactamente estos pasos.
\end{observacion}

% Figura: problema 4.8. Condensador esférico con dos dieléctricos hemisféricos.
\begin{figure}[t]
\centering
\begin{subfigure}[b]{0.5\textwidth}\centering
\begin{tikzpicture}[scale=1.0]
  \fill[dielec] (0,2.1) arc (90:270:2.1) -- cycle;
  \fill[dielecB] (0,-2.1) arc (-90:90:2.1) -- cycle;
  \draw[texto!70!fondo, line width=0.5pt] (0,-2.1) -- (0,2.1);
  \draw[texto, line width=1.1pt] (0,0) circle (2.1);
  \shade[ball color=conductorluz] (0,0) circle (0.85);
  \draw[texto, line width=0.6pt] (0,0) circle (0.85);
  \draw[gauss] (0,0) circle (1.5);
  \node[etiq] at (-0.9,1.62) {$\varepsilon_1$};
  \node[etiq] at (0.9,1.62) {$\varepsilon_2$};
  \draw[texto, line width=0.4pt, -{Stealth[length=1.6mm]}] (0,0) -- (40:0.85);
  \node[font=\small, text=black] at (0.12,0.12) [above left=-1pt] {$a$};
  \draw[acentoB, line width=0.5pt, -{Stealth[length=1.6mm]}] (0,0) -- (-25:1.5);
  \node[etiq, text=acentoB] at (-20:1.25) [below=1pt] {$r$};
  \draw[texto, line width=0.4pt, -{Stealth[length=1.6mm]}] (0,0) -- (-125:2.1);
  \node[etiq] at (-125:1.75) [left=1pt] {$b$};
  \node[etiq] at (-1.55,0.25) {$Q$};
  \draw[texto, line width=0.4pt, -{Stealth[length=1.5mm]}] (-1.38,0.18) -- (165:0.9);
  % tierra
  \draw[texto, line width=0.6pt] (2.1,0) -- (2.75,0) -- (2.75,-0.25);
  \pic at (2.75,-0.25) {tierra};
  \node[etiqt, anchor=west] at (2.82,0.15) {$\phi=0$};
\end{tikzpicture}
\caption{Geometría y superficie gaussiana}
\end{subfigure}\hfill
\begin{subfigure}[b]{0.47\textwidth}\centering
\begin{tikzpicture}[scale=0.42]
  \newcommand{\mini}[1]{\begin{scope}[shift={(0,#1)}]
     \fill[dielec] (0,2.1) arc (90:270:2.1) -- cycle;
     \fill[dielecB] (0,-2.1) arc (-90:90:2.1) -- cycle;
     \draw[texto, line width=0.6pt] (0,0) circle (2.1);
     \fill[fondo] (0,0) circle (0.95);\end{scope}}
  % (i) interfase: sin carga
  \mini{5.6}
  \draw[acentoC, line width=1.6pt] (0,0.95+5.6) -- (0,2.1+5.6) (0,-0.95+5.6) -- (0,-2.1+5.6);
  \draw[texto, line width=0.5pt] (0,5.6) circle (0.95);
  \node[etiq, anchor=west, align=left, font=\footnotesize] at (2.5,5.6) {interfase plana: $E_n=D_n=0$,\\ sin carga libre ni ligada};
  % (ii) carga total uniforme
  \mini{0}
  \draw[acentoA, line width=2.0pt] (0,0) circle (0.95);
  \node[etiq, anchor=west, align=left, font=\footnotesize] at (2.5,0) {carga total en $r=a$:\\ $\sigmaf+\sigmaP=\eo E(a)$, uniforme};
  % (iii) carga libre mayor en el medio de mayor permitividad
  \mini{-5.6}
  \draw[cargapos, line width=1.0pt] (0,0.95-5.6) arc (90:270:0.95);
  \draw[cargapos, line width=2.6pt] (0,-0.95-5.6) arc (-90:90:0.95);
  \node[etiq, anchor=west, align=left, font=\footnotesize] at (2.5,-5.6) {carga libre: $\sigma_{\mathrm{f},i}=\varepsilon_i E(a)$,\\ mayor donde $\varepsilon$ es mayor};
\end{tikzpicture}
\caption{Reparto de la carga en $r=a$ ($\varepsilon_2>\varepsilon_1$)}
\end{subfigure}
\caption{\Cref{prob:dos-hemisferios}. (a) Esfera interior de radio $a$ y carga $Q$; esfera exterior de radio $b$ a tierra; el hueco
está lleno por dos dieléctricos separados por un plano ecuatorial. (b) El campo es radial e igual en ambos
medios, pero la carga libre de la esfera interior no es uniforme: se concentra en el hemisferio del
dieléctrico de mayor permitividad, de modo que la carga total (libre más ligada) sí resulta uniforme.}
\label{fig:p48}
\end{figure}


% -----------------------------------------------------------------------------
\subsection{Problema 4.9. Electrete en un condensador cortocircuitado}\label{sec:p49}

\begin{problema}[Electrete entre armaduras unidas]\label{prob:electrete-condensador}
Una plancha de electrete con polarización $\vP$ uniforme se coloca entre las armaduras de un condensador plano como se muestra en la
\cref{fig:p49}. Calcula los campos dentro y fuera del electrete cuando las armaduras del condensador se unen mediante un cable conductor.
Desprecia el efecto de bordes.
\end{problema}

\begin{resolucion}
\idea La polarización es uniforme e independiente de cualquier campo. El electrete tiene cargas ligadas $\pm P$ en sus caras, que inducen
carga en las armaduras; al estar unidas por el cable, sus potenciales son iguales. Sea $a$ el espesor del electrete y $b$ el del hueco de aire, y
$\uv z$ el sentido de $\vP$, normal a las placas (método M5).

\paso{Paso 1: $\vD$ es el mismo en ambas capas} Por la simetría plana, los campos son normales a las placas y solo dependen de $z$. En cada capa
no hay carga libre y $\dd D/\dd z=0$; en la interfase tampoco, y $D_n$ es continua. Luego $\vD=D\,\uv z$ con $D$ constante en todo el hueco.

\paso{Paso 2: los campos en función de $D$}
\[
  \eo\vEext=\vD\quad(\text{aire}),\qquad \eo\vEint=\vD-\vP\quad(\text{electrete}).
\]
\paso{Paso 3: la condición del cable} Las armaduras están al mismo potencial:
\[
  \phi_{\text{sup}}-\phi_{\text{inf}}=-\int_0^{a+b}E_z\,\dd z=-\big(\Eint\,a+\Eext\,b\big)=0 ,
\]
y sustituyendo los campos del paso~2,
\[
  (D-P)\,a+D\,b=0\quad\Longrightarrow\quad D=\frac{a}{a+b}\,P .
\]
\paso{Paso 4: resultado}
\begin{equation}\label{eq:p49}
  \recuadro{\vEint=-\frac{b}{\eo(a+b)}\,\vP,\qquad \vEext=\frac{a}{\eo(a+b)}\,\vP,\qquad \vD=\frac{a}{a+b}\,\vP .}
\end{equation}
Las cargas inducidas en las armaduras son $\sigma_{\text{inf}}=+D=\frac{a}{a+b}P$ en la inferior (junto a la cara negativa del electrete) y
$\sigma_{\text{sup}}=-D$ en la superior.
\end{resolucion}

\begin{observacion}[Comprobación por superposición de planos]\label{obs:p49}
Hay cuatro planos cargados: $+s$ (armadura inferior), $-P$ y $+P$ (caras del electrete) y $-s$ (armadura superior), con $s=\frac{a}{a+b}P$. Cada uno crea
$\sigma/2\eo$ a cada lado, alejándose si es positivo. En el aire, los tres planos inferiores suman $s/2\eo$ hacia arriba y la armadura superior, que
atrae, otros $s/2\eo$: $\Eext=s/\eo=aP/\eo(a+b)$. En el electrete, $\frac{s-P}{2\eo}-\frac{P}{2\eo}+\frac{s}{2\eo}=\frac{s-P}{\eo}=-\frac{bP}{\eo(a+b)}$. Coincide con \eqref{eq:p49}, como en
la resolución de los seminarios. Sin el cable ($s=0$) se recupera la lámina polarizada aislada, $\Eext=0$ y $\Eint=-P/\eo$; con el hueco cerrado ($b\to0$),
$\Eint\to0$: el electrete queda cortocircuitado por las placas.
\end{observacion}

% Figura: problema 4.9. Lámina de electrete entre las armaduras de un condensador
% unidas por un cable; perfiles de E, D, P y phi (a = 1, b = 2).
\begin{figure}[t]
\centering
\begin{subfigure}[b]{0.52\textwidth}\centering
\begin{tikzpicture}[scale=1.0]
  \def\a{1.0} \def\b{2.0} \def\L{5.0}
  \fill[dielec] (0,0) rectangle (\L,\a);
  \draw[borde] (0,0) rectangle (\L,\a);
  \fill[conductor] (0,-0.14) rectangle (\L,0);
  \fill[conductor] (0,\a+\b) rectangle (\L,\a+\b+0.14);
  \draw[texto, line width=0.5pt] (0,-0.14) rectangle (\L,0) (0,\a+\b) rectangle (\L,\a+\b+0.14);
  % cable
  \draw[texto, line width=0.7pt, rounded corners=6pt] (\L,\a+\b+0.07) -- (\L+0.7,\a+\b+0.07) -- (\L+0.7,-0.07) -- (\L,-0.07);
  % cargas
  \foreach \x in {0.3,0.95,...,4.4} {
    \node[cpos, font=\tiny\bfseries] at (\x,-0.07) {$+$};
    \node[cneg, font=\tiny\bfseries] at (\x,\a+\b+0.07) {$-$};
    \node[cpos, font=\tiny\bfseries] at (\x+0.32,\a-0.13) {$+$};
    \node[cneg, font=\tiny\bfseries] at (\x+0.32,0.13) {$-$};
  }
  % vectores
  \draw[vecP] (0.65,0.22) -- (0.65,0.8); \node[etiqm, font=\footnotesize, anchor=west] at (0.7,0.5) {$\vP$};
  \draw[vecE] (2.2,0.8) -- (2.2,0.32); \node[etiqm, font=\footnotesize, anchor=west] at (2.25,0.52) {$\vEint$};
  \draw[vecE] (2.2,\a+0.4) -- (2.2,\a+0.4+0.6); \node[etiq, font=\footnotesize, anchor=west] at (2.25,\a+0.75) {$\vEext$};
  \draw[vecD] (3.7,\a+0.4) -- (3.7,\a+1.0); \node[etiq, font=\footnotesize, anchor=west] at (3.75,\a+0.75) {$\vD$};
  \draw[vecD] (3.7,0.25) -- (3.7,0.85); \node[etiqm, font=\footnotesize, anchor=west] at (3.75,0.55) {$\vD$};
  % cotas
  \draw[cota] (-0.35,0) -- (-0.35,\a) node[midway, left, etiq] {$a$};
  \draw[cota] (-0.35,\a) -- (-0.35,\a+\b) node[midway, left, etiq] {$b$};
  \node[etiqt, anchor=south] at (\L/2,\a+\b+0.16) {$\sigma_{\text{sup}}=-\tfrac{a}{a+b}P$};
  \node[etiqt, anchor=north] at (\L/2,-0.17) {$\sigma_{\text{inf}}=+\tfrac{a}{a+b}P$};
\end{tikzpicture}
\caption{Geometría y cargas}
\end{subfigure}\hfill
\begin{subfigure}[b]{0.45\textwidth}\centering
\begin{tikzpicture}
\begin{axis}[width=6.3cm, height=5.6cm, xmin=-0.9, xmax=1.15, ymin=0, ymax=3,
  ytick={0,1,3}, yticklabels={$0$,$a$,$a+b$}, ylabel={altura}, xtick={-0.6667,0,0.3333,0.6667,1},
  xticklabels={$-\frac23$,$0$,$\frac13$,$\frac23$,$1$}, grid=major,
  legend style={at={(0.5,1.03)}, anchor=south, legend columns=2, font=\tiny, /tikz/every even column/.append style={column sep=3pt}}]
  \fill[dielec, opacity=0.6] (axis cs:-0.9,0) rectangle (axis cs:1.15,1);
  \addplot[curvaP] coordinates {(1,0) (1,1) (0,1) (0,3)}; \addlegendentry{$P$}
  \addplot[curvaD, line width=2.8pt, opacity=0.45] coordinates {(0.3333,0) (0.3333,3)}; \addlegendentry{$D$}
  \addplot[curvaE] coordinates {(-0.6667,0) (-0.6667,1)}; \addlegendentry{$\eo E$}
  \addplot[curvaV] coordinates {(0,0) (0.6667,1) (0,3)}; \addlegendentry{$\phi\,\eo/(Pa)$}
  \addplot[curvaE, forget plot] coordinates {(0.3333,1) (0.3333,3)};
  \addplot[curvaE, forget plot, dotted, line width=0.5pt] coordinates {(-0.6667,1) (0.3333,1)};
\end{axis}
\end{tikzpicture}
\caption{Perfiles ($b=2a$, unidades de $P$)}
\end{subfigure}
\caption{\Cref{prob:electrete-condensador}. (a) Electrete de espesor $a$ y polarización uniforme $\vP$ sobre la armadura inferior; las
armaduras están unidas por un cable, de modo que $\phi_{\text{sup}}=\phi_{\text{inf}}$. Sin carga libre en la interfase, $\vD$ es
el mismo en el electrete y en el aire, y la condición $\oint\vE\cdot\dl=0$ obliga a que $\vE$ cambie de sentido.
(b) Con $b=2a$: $\eo\Eint=-\tfrac23P$, $\eo\Eext=\tfrac13P$, $D=\tfrac13P$; el potencial sube en el electrete y baja en el aire hasta
volver al valor inicial. En el aire $D=\eo\Eext$, y ambas curvas coinciden.}
\label{fig:p49}
\end{figure}


% -----------------------------------------------------------------------------
\subsection{Problema 4.10. Dieléctricos en paralelo}\label{sec:p410}

\begin{problema}[Dos dieléctricos en paralelo]\label{prob:paralelo}
Entre las placas paralelas de un condensador plano al que aplicamos una diferencia de potencial $V_0$ se introducen dos dieléctricos de
permitividades relativas $\varepsilon_{r1}$ y $\varepsilon_{r2}$ como se indica en la \cref{fig:p410-412}\,(a). Calcula la carga en cada placa y la capacidad del
sistema. Desprecia el efecto de bordes.
\end{problema}

\begin{resolucion}
\idea La interfase entre los dieléctricos es perpendicular a las placas, de modo que los campos son tangentes a ella (método M4, paralelo).

\paso{Paso 1: el campo es común} La componente tangencial de $\vE$ es continua, $E_1=E_2\equiv E$. Como las placas son equipotenciales y separadas $d$,
\[
  E_1=E_2=E=\frac{V_0}{d}.
\]
\paso{Paso 2: carga en cada zona} En la superficie de la placa, $\sigmaf=D_n$: $\sigma_i=D_i=\varepsilon_{ri}\eo E=\varepsilon_{ri}\eo V_0/d$, y con las áreas $A_1$ y $A_2$,
\[
  \recuadro{Q_1=V_0\,\frac{A_1\varepsilon_{r1}\eo}{d},\qquad Q_2=V_0\,\frac{A_2\varepsilon_{r2}\eo}{d},}
\]
y la capacidad es la suma de las de ambos tramos,
\[
  \recuadro{C=\frac{Q_1+Q_2}{V_0}=\frac{A_1\varepsilon_{r1}\eo}{d}+\frac{A_2\varepsilon_{r2}\eo}{d}=C_1+C_2 .}
\]
\end{resolucion}

\begin{observacion}[Paralelo]\label{obs:p410}
Los dos tramos comparten la diferencia de potencial y se reparten la carga: son dos condensadores en paralelo. La solución de campo
uniforme cumple todas las condiciones (placas equipotenciales, $\vE_t$ continuo y $D_n$ continua, trivialmente, en la interfase), y es única. El
mismo argumento resolvió el \cref{prob:dos-hemisferios}.
\end{observacion}

% -----------------------------------------------------------------------------
\subsection{Problema 4.11. Dieléctrico en serie con aire}\label{sec:p411}

\begin{problema}[Condensador parcialmente lleno]\label{prob:serie}
Considera un condensador de placas paralelas de área $A$ y separación $d$. Entre sus placas se aplica una diferencia de potencial $V_0$ y se
introduce un dieléctrico de permitividad $\varepsilon$ que llena parcialmente el espacio entre placas, como se indica en la \cref{fig:p410-412}\,(b).
Calcula los campos dentro y fuera del dieléctrico y la capacidad del sistema. Desprecia el efecto de bordes.
\end{problema}

\begin{resolucion}
\idea Ahora la interfase es paralela a las placas y los campos son normales a ella (método M4, serie). Sean $a$ el espesor del dieléctrico y $b=d-a$ el
del aire.

\paso{Paso 1: $\vD$ es común} Sin carga libre en la interfase, $D_n$ es continua: $\eo\Eext=\varepsilon\Eint\equiv D$.

\paso{Paso 2: la batería} $V_0=\Eint\,a+\Eext\,b=\Eint\Big(a+\dfrac{\varepsilon}{\eo}\,b\Big)=\dfrac{\Eint}{\eo}\big(\eo a+\varepsilon b\big)$, de donde
\[
  \recuadro{\Eint=V_0\,\frac{\eo}{a\eo+b\varepsilon},\qquad \Eext=V_0\,\frac{\varepsilon}{a\eo+b\varepsilon},\qquad
  C=\frac{DA}{V_0}=A\,\frac{\varepsilon\eo}{a\eo+b\varepsilon} .}
\]
\paso{Paso 3: comprobación como serie} Con $C_1=\eo A/b$ (aire) y $C_2=\varepsilon A/a$ (dieléctrico),
\[
  \frac{C_1C_2}{C_1+C_2}=\frac{A^2\eo\varepsilon/(ab)}{A(\eo a+\varepsilon b)/(ab)}=A\,\frac{\varepsilon\eo}{a\eo+b\varepsilon}. \qedhere
\]
\end{resolucion}

\begin{observacion}[Límites]\label{obs:p411}
Si $\varepsilon\to\eo$, $C\to\eo A/(a+b)=\eo A/d$. Si $\varepsilon\to\infty$, $\Eint\to0$ y $C\to\eo A/b$: el dieléctrico se comporta como un conductor y la separación efectiva se
reduce al hueco de aire, que es el \cref{prob:lamina}. Casi toda la caída de potencial se produce en el aire, donde el campo es $\er$ veces mayor
(\cref{fig:p411-perfil}\,(a)): por eso una capa de aire en serie con un aislante es el punto débil frente a la ruptura dieléctrica.
\end{observacion}

% -----------------------------------------------------------------------------
\subsection{Problema 4.12. Lámina conductora intercalada}\label{sec:p412}

\begin{problema}[Lámina conductora]\label{prob:lamina}
Entre las placas de un condensador plano-paralelo de área $A$ y separación $d$ se introduce una lámina conductora de grosor $e$ como se indica en la
\cref{fig:p410-412}\,(c). ¿Cómo varía la capacidad del sistema?
\end{problema}

\begin{resolucion}
\idea Dentro del conductor $\vE=\vcero$; en sus caras se inducen cargas $\mp Q$ que dejan inalterado el campo de los huecos (método M4, serie).

\paso{Paso 1: campos} Con cargas $\pm Q$ en las placas, la lámina adquiere $-Q$ en la cara frente a la placa positiva y $+Q$ en la otra (regla de las caras del
Tema~3~\cite{T3}). En cada hueco, $D=\sigma=Q/A$ y $E=\sigma/\eo$; en el metal, $E=0$.

\paso{Paso 2: diferencia de potencial} Con $a$ y $b$ los espesores de los huecos, $a+b=d-e$:
\[
  \Delta\phi=\frac{\sigma}{\eo}\,a+0\cdot e+\frac{\sigma}{\eo}\,b=\frac{\sigma}{\eo}(d-e)
  \quad\Longrightarrow\quad
  \recuadro{C'=\frac{Q}{\Delta\phi}=\frac{\eo A}{d-e}>\frac{\eo A}{d}=C .}
\]
\end{resolucion}

\begin{observacion}[Independencia de la posición]\label{obs:p412}
La capacidad aumenta y no depende de dónde se coloque la lámina, solo de su espesor: equivale a dos condensadores en serie, $\eo A/a$ y $\eo A/b$, cuya
serie solo depende de $a+b$. Es el límite $\varepsilon\to\infty$ del \cref{prob:serie} (\cref{fig:p411-perfil}\,(b)).
\end{observacion}

% Figura: problemas 4.10, 4.11 y 4.12. Condensadores planos con dieléctricos
% «en paralelo», «en serie» y con una lámina conductora intercalada.
\begin{figure}[t]
\centering
\begin{subfigure}[b]{0.3\textwidth}\centering
\begin{tikzpicture}[scale=0.8]
  \fill[dielecB] (0,1.7) rectangle (2.0,3.4);
  \fill[dielec] (0,0) rectangle (2.0,1.7);
  \draw[texto!70!fondo, line width=0.5pt] (0,1.7) -- (2.0,1.7);
  \draw[placa] (0,-0.1) -- (0,3.5);
  \draw[placa] (2.0,-0.1) -- (2.0,3.5);
  \draw[cargapos, line width=1.6pt] (-0.06,0) -- (-0.06,3.4);
  \node[etiq] at (1.0,3.05) {$\varepsilon_{r1}$};
  \node[etiq] at (1.0,1.35) {$\varepsilon_{r2}$};
  \draw[vecE] (0.45,2.35) -- (1.55,2.35);
  \draw[vecE] (0.45,0.65) -- (1.55,0.65);
  \node[etiq, anchor=east] at (-0.15,2.95) {$A_1$};
  \node[etiq, anchor=east] at (-0.15,2.3) {$Q_1$};
  \node[etiq, anchor=east] at (-0.15,1.25) {$A_2$};
  \node[etiq, anchor=east] at (-0.15,0.6) {$Q_2$};
  \draw[cota] (0,3.8) -- (2.0,3.8) node[midway, above, etiq] {$d$};
  % batería
  \draw[texto, line width=0.6pt] (0,-0.1) -- (0,-0.95) -- (0.8,-0.95);
  \draw[texto, line width=0.6pt] (1.2,-0.95) -- (2.0,-0.95) -- (2.0,-0.1);
  \draw[texto, line width=1.1pt] (0.85,-0.72) -- (0.85,-1.18);
  \draw[texto, line width=0.6pt] (1.12,-0.82) -- (1.12,-1.08);
  \node[etiq] at (1.0,-1.45) {$V_0$};
\end{tikzpicture}
\caption{Problema 4.10}
\end{subfigure}\hfill
\begin{subfigure}[b]{0.34\textwidth}\centering
\begin{tikzpicture}[scale=0.8]
  \fill[dielec] (0,0) rectangle (3.2,1.2);
  \draw[texto!70!fondo, line width=0.5pt] (0,1.2) -- (3.2,1.2);
  \draw[placa] (-0.1,0) -- (3.3,0);
  \draw[placa] (-0.1,2.6) -- (3.3,2.6);
  \node[etiq] at (2.95,0.6) {$\varepsilon$};
  \node[etiq] at (2.95,1.9) {$\eo$};
  \draw[vecD] (0.6,0.15) -- (0.6,1.05); \draw[vecD] (0.6,1.45) -- (0.6,2.35);
  \node[etiq, anchor=west, font=\footnotesize] at (0.68,1.9) {$\vD$};
  \draw[vecE] (1.6,0.4) -- (1.6,0.8); \draw[vecE] (1.6,1.35) -- (1.6,2.45);
  \node[etiq, anchor=west, font=\footnotesize] at (1.68,1.9) {$\vEext$};
  \node[etiqm, anchor=west, font=\footnotesize] at (1.68,0.6) {$\vEint$};
  \draw[cota] (-0.35,0) -- (-0.35,1.2) node[midway, left, etiq] {$a$};
  \draw[cota] (-0.35,1.2) -- (-0.35,2.6) node[midway, left, etiq] {$b$};
  % batería a la derecha
  \draw[texto, line width=0.6pt] (3.3,2.6) -- (3.95,2.6) -- (3.95,1.5);
  \draw[texto, line width=0.6pt] (3.95,1.1) -- (3.95,0) -- (3.3,0);
  \draw[texto, line width=1.1pt] (3.72,1.45) -- (4.18,1.45);
  \draw[texto, line width=0.6pt] (3.82,1.18) -- (4.08,1.18);
  \node[etiq, anchor=west] at (4.2,1.3) {$V_0$};
\end{tikzpicture}
\caption{Problema 4.11}
\end{subfigure}\hfill
\begin{subfigure}[b]{0.32\textwidth}\centering
\begin{tikzpicture}[scale=0.8]
  \fill[conductor] (0,1.2) rectangle (3.2,1.65);
  \draw[texto, line width=0.5pt] (0,1.2) rectangle (3.2,1.65);
  \draw[placa] (-0.1,0) -- (3.3,0);
  \draw[placa] (-0.1,3.0) -- (3.3,3.0);
  \node[etiq] at (0.45,0.6) {$\eo$};
  \node[etiq] at (0.45,2.33) {$\eo$};
  \node[etiq, anchor=west] at (3.35,3.0) {$+Q$};
  \node[etiq, anchor=west] at (3.35,1.73) {$-Q$};
  \node[etiq, anchor=west] at (3.35,1.12) {$+Q$};
  \node[etiq, anchor=west] at (3.35,0) {$-Q$};
  \draw[vecE] (1.2,2.85) -- (1.2,1.8); \draw[vecE] (1.2,1.05) -- (1.2,0.15);
  \node[etiq, anchor=west, font=\footnotesize] at (1.28,2.33) {$\sigma/\eo$};
  \node[etiq, anchor=west, font=\footnotesize] at (1.28,0.6) {$\sigma/\eo$};
  \node[font=\footnotesize, text=black] at (1.9,1.425) {$\vE=\vcero$};
  \draw[cota] (-0.35,0) -- (-0.35,1.2) node[midway, left, etiq] {$a$};
  \draw[cota] (-0.35,1.2) -- (-0.35,1.65) node[midway, left, etiq] {$e$};
  \draw[cota] (-0.35,1.65) -- (-0.35,3.0) node[midway, left, etiq] {$b$};
\end{tikzpicture}
\caption{Problema 4.12}
\end{subfigure}
\caption{Condensadores planos de área $A$ y separación $d$ parcialmente llenos. (a) Dos dieléctricos con la interfase
perpendicular a las placas: el campo es el mismo en ambos ($V_0/d$), pero la carga libre se reparte en
proporción a $\varepsilon_{ri}A_i$; el sistema equivale a dos condensadores en paralelo. (b) Interfase paralela a las
placas: $\vD$ es el mismo en ambas capas y el campo es menor en el dieléctrico; equivale a dos condensadores en
serie. (c) Una lámina conductora de espesor $e$ anula el campo en su interior y reduce la separación efectiva
a $d-e$, cualquiera que sea su posición.}
\label{fig:p410-412}
\end{figure}

% Figura: problemas 4.11 y 4.12. Perfiles de campo y potencial frente a la altura.
\begin{figure}[t]
\centering
\begin{subfigure}[b]{0.48\textwidth}\centering
\begin{tikzpicture}
\begin{axis}[width=6.6cm, height=5.4cm, xmin=0, xmax=1.65, ymin=0, ymax=2,
  ytick={0,1,2}, yticklabels={$0$,$a$,$d$}, ylabel={altura}, xtick={0,0.5,1,1.5},
  xticklabels={$0$,$\frac12$,$1$,$\frac32$}, grid=major,
  legend style={at={(0.5,1.03)}, anchor=south, legend columns=3, font=\tiny, /tikz/every even column/.append style={column sep=3pt}}]
  \fill[dielec, opacity=0.6] (axis cs:0,0) rectangle (axis cs:1.65,1);
  \addplot[curvaD, line width=2.8pt, opacity=0.45] coordinates {(1.5,0) (1.5,2)}; \addlegendentry{$D\,d/(\eo V_0)$}
  \addplot[curvaE] coordinates {(0.5,0) (0.5,1)}; \addlegendentry{$E\,d/V_0$}
  \addplot[curvaV] coordinates {(0,0) (0.25,1) (1,2)}; \addlegendentry{$\phi/V_0$}
  \addplot[curvaE, forget plot] coordinates {(1.5,1) (1.5,2)};
  \addplot[curvaE, forget plot, dotted, line width=0.5pt] coordinates {(0.5,1) (1.5,1)};
  \node[etiqt, anchor=north east] at (axis cs:1.35,0.95) {$\varepsilon=3\eo$};
\end{axis}
\end{tikzpicture}
\caption{Problema 4.11 ($a=b$, $\er=3$)}
\end{subfigure}\hfill
\begin{subfigure}[b]{0.48\textwidth}\centering
\begin{tikzpicture}
\begin{axis}[width=6.6cm, height=5.4cm, xmin=0, xmax=1.4, ymin=0, ymax=3,
  ytick={0,1,1.6,3}, yticklabels={$0$,$a$,$a+e$,$d$}, ylabel={altura}, xtick={0,0.5,1},
  xticklabels={$0$,$\frac12$,$1$}, grid=major,
  legend style={at={(0.5,1.03)}, anchor=south, legend columns=3, font=\tiny, /tikz/every even column/.append style={column sep=3pt}}]
  \fill[conductor, opacity=0.45] (axis cs:0,1) rectangle (axis cs:1.4,1.6);
  \addplot[curvaE] coordinates {(1,0) (1,1) (0,1) (0,1.6) (1,1.6) (1,3)}; \addlegendentry{$E\,(d-e)/V_0$}
  \addplot[curvaV] coordinates {(0,0) (0.4167,1) (0.4167,1.6) (1,3)}; \addlegendentry{$\phi/V_0$}
  \node[etiqt] at (axis cs:0.95,1.3) {metal};
\end{axis}
\end{tikzpicture}
\caption{Problema 4.12 ($a=1$, $e=0{,}6$, $b=1{,}4$)}
\end{subfigure}
\caption{Perfiles en función de la altura sobre la placa inferior (que se toma a potencial nulo). (a) Capa dieléctrica de
espesor $a$ bajo una capa de aire de espesor $b$: $D$ es continuo, el campo salta en el factor $\er$ y el potencial
cae sobre todo en el aire. (b) Lámina conductora: el campo es nulo dentro del metal y el potencial es allí
constante; fuera, el campo es el de un condensador de separación $d-e$.}
\label{fig:p411-perfil}
\end{figure}


% -----------------------------------------------------------------------------
\subsection{Resumen de resultados y lista de control}\label{sec:resumen-seminarios}

\begin{table}[H]
\centering
\caption{Resultados de los seminarios del Tema 4.}\label{tab:resultados}
\small
\renewcommand{\arraystretch}{1.55}
\begin{tabularx}{\textwidth}{@{}c>{\raggedright\arraybackslash}Xc@{}}
\toprule
Problema & Resultado & Método\\
\midrule
4.1 & $\phi=\dfrac{Pa^3}{3\eo}\dfrac{z}{\abs z^3}$, $E_z=\dfrac{2Pa^3}{3\eo\abs z^3}$ ($\abs z>a$);\quad $\phi=\dfrac{Pz}{3\eo}$, $E_z=-\dfrac{P}{3\eo}$ ($\abs z<a$) & M1\\
4.2 & $\vD=\dfrac{Q}{4\pi r^2}\uv r$,\quad $\vE=\dfrac{Q}{4\pi\varepsilon r^2}\uv r$,\quad $\vP=(\varepsilon-\eo)\dfrac{Q}{4\pi\varepsilon r^2}\uv r$ & M3\\
4.3 & $\vD=\sigma\uv z$,\quad $\vE=\dfrac{\sigma}{\varepsilon}\uv z$,\quad $C=\dfrac{\varepsilon A}{d}=\er C_0$ & M3\\
4.4 & $\vE=\dfrac{\vP}{2\eo}(\cos\theta_2-\cos\theta_1)$ fuera, $-\dfrac{\vP}{2\eo}(2+\cos\theta_1-\cos\theta_2)$ dentro;\quad $\vD=\dfrac{\vP}{2}(\cos\theta_2-\cos\theta_1)$ & M1\\
4.5 & $\rhoP=-3A$,\quad $\sigmaP=A\ell/2$,\quad $\QP=0$ & M1\\
4.6 & $\vD=\vcero$;\quad $\vE=-\dfrac{k}{\eo r^2}\uv r$ en $a<r<2a$, nulo fuera & M2\\
4.7 & $\vD=\vcero$;\quad $\vE=-\dfrac{kr}{2\eo}\uv r$ en $a<r<b$, nulo fuera;\quad $\rhoP=-k$ & M2\\
4.8 & $\phi=\dfrac{Q}{2\pi(\varepsilon_1+\varepsilon_2)}\Big(\dfrac1r-\dfrac1b\Big)$,\quad $C=2\pi(\varepsilon_1+\varepsilon_2)\dfrac{ab}{b-a}$,\quad $\sigma_i=\dfrac{\varepsilon_iQ}{2\pi a^2(\varepsilon_1+\varepsilon_2)}$ & M3, M4\\
4.9 & $\vEint=-\dfrac{b\,\vP}{\eo(a+b)}$,\quad $\vEext=\dfrac{a\,\vP}{\eo(a+b)}$,\quad $\vD=\dfrac{a\,\vP}{a+b}$ & M5\\
4.10 & $Q_i=V_0\dfrac{A_i\varepsilon_{ri}\eo}{d}$,\quad $C=\dfrac{A_1\varepsilon_{r1}\eo}{d}+\dfrac{A_2\varepsilon_{r2}\eo}{d}$ & M4\\
4.11 & $\Eint=\dfrac{V_0\eo}{a\eo+b\varepsilon}$,\quad $\Eext=\dfrac{V_0\varepsilon}{a\eo+b\varepsilon}$,\quad $C=\dfrac{A\varepsilon\eo}{a\eo+b\varepsilon}$ & M4\\
4.12 & $C'=\dfrac{\eo A}{d-e}$ (independiente de la posición) & M4\\
\bottomrule
\end{tabularx}
\end{table}

\begin{metodo}{Lista de control}
\begin{itemize}
  \item ¿He distinguido carga libre (la que controlo) de carga ligada (la que dispone el material)? ¿Uso Gauss para $\vD$ solo con la libre?
  \item Antes de afirmar $\vD=\vcero$, ¿he comprobado que no hay carga libre \emph{y} que $\vP$ es irrotacional, también en las superficies?
  \item ¿He tomado en cada superficie la normal \emph{exterior al dieléctrico} al calcular $\sigmaP=\vP\cdot\nn$?
  \item ¿Se anula la carga ligada total? ¿Coincide el campo con el que da Gauss para $\vE$ con todas las cargas?
  \item ¿Uso $\oiint\vE\cdot\dS=\qfenc/\varepsilon$ solo si un único medio \LHI{} llena todo el espacio relevante?
  \item En las interfases sin carga libre, ¿he impuesto $D_n$ continua y $\vE_t$ continuo (y no al revés)?
  \item Condensadores: interfase paralela a las placas, $\vD$ común y serie; perpendicular, $\vE$ común y paralelo.
  \item ¿Recupera el resultado los límites conocidos: $\varepsilon\to\eo$ (vacío), $\varepsilon\to\infty$ (conductor), $b\to0$, $a\gg\ell$, $a\ll\ell$?
\end{itemize}
\end{metodo}
